\documentclass[12pt,letterpaper]{article}
\usepackage[utf8]{inputenc}
\usepackage[spanish]{babel}
\usepackage{amsmath}
\usepackage{amsfonts}
\usepackage{amssymb}
\usepackage[margin=2.5cm]{geometry}

\begin{document}

\begin{center}
    \Large \textbf{Evaluación de Examen 2: Unidad II} \\
    \large Física para Ciencias de la Salud (005-1713) \\
    \vspace{0.5cm}
    \textbf{Estudiante:} Ybranmaris Maribao \quad \textbf{C.I.:} 36.750.392
    \vspace{0.3cm} \\
    \large \textbf{Calificación Final: 10/10 puntos}
\end{center}

\vspace{0.5cm}

\section*{Análisis de Resultados}

\subsection*{1. Cinemática (Aceleración media)}
\textbf{Procedimiento y resultado:} Extrae de manera estructurada los datos iniciales y plantea la ecuación de aceleración media ($a = \frac{v_f - v_i}{\Delta t}$). Realiza la sustitución correcta restando las velocidades y dividiendo entre el tiempo ($0,6 \text{ m/s} / 0,15 \text{ s}$). El cálculo aritmético arroja el resultado exacto de $4 \text{ m/s}^2$ con un impecable manejo dimensional. \\
\textbf{Veredicto:} \textbf{Correcto} (2/2 pts).

\subsection*{2. Dinámica (Leyes de Newton)}
\textbf{Procedimiento y resultado:} Muestra un despeje impecable de la Segunda Ley de Newton para hallar la aceleración ($a = \frac{F_{neta}}{m}$). Sustituye los valores dividiendo $150 \text{ N}$ entre $25 \text{ kg}$, logrando un resultado perfecto de $6 \text{ m/s}^2$. \\
\textbf{Veredicto:} \textbf{Correcto} (2/2 pts).

\subsection*{3. Trabajo Mecánico}
\textbf{Procedimiento y resultado:} Extrae los datos e identifica de forma sobresaliente la fórmula general del trabajo al incluir explícitamente la función trigonométrica ($W = f \cdot d \cdot \cos(0)$). Efectúa la multiplicación reconociendo que el coseno es igual a 1 ($400 \text{ N} \cdot 0,8 \text{ m} \cdot 1$). El cálculo da exactamente $320 \text{ J}$. \\
\textbf{Veredicto:} \textbf{Correcto} (2/2 pts).

\subsection*{4. Energía Potencial}
\textbf{Procedimiento y resultado:} Extrae rigurosamente los datos del problema (masa, altura y aceleración de gravedad). Plantea el producto de las tres variables según la ecuación de la energía potencial gravitatoria ($E_p = m \cdot g \cdot h$) integrando las unidades dentro de la sustitución ($1,2 \text{ kg} \cdot 9,8 \text{ m/s}^2 \cdot 1,5 \text{ m}$). Obtiene el resultado verificado y exacto de $17,64 \text{ J}$. \\
\textbf{Veredicto:} \textbf{Correcto} (2/2 pts).

\subsection*{5. Potencia Mecánica}
\textbf{Procedimiento y resultado:} Extrae los datos y relaciona el trabajo efectuado y el tiempo planteando la fracción respectiva ($P = \frac{W}{\Delta T}$). El desarrollo de la división de $75 \text{ J} / 60 \text{ s}$ arroja la cifra correcta y verificada de $1,25 \text{ W}$. \\
\textbf{Veredicto:} \textbf{Correcto} (2/2 pts).

\vspace{0.5cm}
\section*{Comentarios Generales y Sugerencias}
¡Muchas felicidades por obtener una calificación perfecta! El examen demuestra un nivel de excelencia muy alto, con una sólida asimilación de los conceptos teóricos y una gran habilidad matemática.
\begin{itemize}
    \item El formato de resolución utilizado (separando de forma impecable en \textit{Datos}, \textit{Fórmula} y luego el desarrollo algebraico) es el estándar de calidad esperado en reportes de laboratorio e investigación clínica.
    \item Se valora enormemente la rigurosidad formal, como el uso del delta ($\Delta t$) para representar el intervalo de tiempo, y la inclusión completa de la función trigonométrica en el cálculo del Trabajo Mecánico.
    \item ¡Sigue con este mismo entusiasmo y nivel de excelencia en la próxima unidad!
\end{itemize}

\end{document}